 \documentclass[10pt,a4paper]{article}
\newcommand{\sharedmacros}{../../../SharedMacros}
\include{\sharedmacros/ccategories.macros}
\include{\sharedmacros/ermacros}
\include{\sharedmacros/erdiagram}
\include{\sharedmacros/theorems.macros}
\include{\sharedmacros/general.macros.v2}

\newcommand{\implicitref}[1]{%
  \index[lemmas]{#1!implicit use}%
}

\include{\sharedmacros/diagram.macros}
\include{\sharedmacros/rangeCategories.macros.tex}
\usepackage{xparse}
\usepackage{xstring} % for some diagram macros
\usepackage{fancyvrb} % used for documenting examples of macros 
                      % along with verbatim text of example.
\usepackage{anyfontsize}
\usepackage{mathtools}  % why have we got this?
\usepackage{mdframed}
\usepackage{alltt}    
\usepackage{cmll}
\usepackage{ulem}
\renewcommand{\ttdefault}{txtt}
\usepackage[left=1.5cm, right=4cm, marginparwidth=3cm, top=2cm, bottom=2.0cm]{geometry}
\usepackage{framed}  %used for box framing for example in asides
\usepackage[font=small]{caption}
\usepackage{MnSymbol}
\setlength{\captionmargin}{2cm}
\theoremstyle{remark}
\newtheorem*{lemma*}{Lemma}
\renewcommand{\term}[1]{\textit{#1}}  %SIMPLE UNINDEXED VERSION
\usepackage{imakeidx}
\makeindex[name=definitions, title=Index of Definitions]
\makeindex[name=lemmas, title=Index of Lemmas]
\usepackage{makeidx}
\makeindex

\newcommand{\fgsourcediag}{$\binarysourcediag{a}{b}{c}{f}{g}$}
\newcommand{\fgsinkdiag}{$\binarysinkdiag{a}{b}{c}{f}{g}$}
\newcommand{\FGsourcediag}{$\binarysourcediag{a}{b}{c}{F}{G}$}
\newcommand{\FGsinkdiag}{$\binarysinkdiag{a}{b}{c}{F}{G}$}

\newcommand{\equsing}[1]{\overset{\text{#1}}{=}}

%%%%%%%%%%%%%%%%%%%%%%%%%%%%%%%%%%
\newcommand{\srcFolder}{src}
%
%\includeonly{\srcFolder/MToD.rangeCategories.zigzagcommand.tests.tex}
%\includeonly{\srcFolder/MToD.rangeCategories.introduction.tex}
%
%\includeonly{\srcFolder/MToD.rangeCategories.supportCategories.tex}
%\includeonly{\srcFolder/MToD.rangeCategories.cosupportCategories.tex}
\includeonly{\srcFolder/MToD.rangeCategories.bisupportCategories.tex}
%\includeonly{\srcFolder/MToD.rangeCategories.posetEnrichedBisupportCategories.tex}
%\includeonly{\srcFolder/MToD.rangeCategories.zigzags.tex}
%\includeonly{\srcFolder/MToD.rangeCategories.bisupportWithConverse.tex}
%\includeonly{\srcFolder/MToD.rangeCategories.restrictions_and_support.tex}
%\includeonly{\srcFolder/MToD.rangeCategories.ranges_sections_and_weak_range.tex} 
%\includeonly{\srcFolder/MToD.rangeCategories.shortening.tex}
%\includeonly{\srcFolder/MToD.rangeCategories.carry.tex}
%\includeonly{\srcFolder/MToD.rangeCategories.Rel_category.tex}
%\section{Goodness of Data Specifications}
%\includeonly{\srcFolder/MToD.rangeCategories.goodness_principles.tex}
%\includeonly{\srcFolder/MToD.rangeCategories.goodness_criteria.tex}
%\section{Goodness Lemmas}
%\includeonly{\srcFolder/MToD.rangeCategories.lemma2A.path_equivalence_representation.tex}
%\includeonly{\srcFolder/MToD.rangeCategories.lemma2B.fd_representation.tex}
%\includeonly{\srcFolder/MToD.rangeCategories.lemma2C.referential_dependency_representation.tex}
%%%%%%%%%%%%%%%%%%%%%%%%%%%%%%%%%%%%

\begin{document}
\title{Range Categories and the Theory of Data\\[0.5em]
\large Foundations of the Mathematical Theory of Data}

% abstract here for ams
\author{John Cartmell}
\date{}
\maketitle
\begin{aside}
14th October 2026. I have gravitated toward bisupport catgories because 
many of the proofs regarding tightness and the establishing of a category of tight zigzags seem to befit from appeals to symmetry. It is mathematically expedient to symmetry of support and cosupport. I became unstuck however because if the the partial order defined in the restriction category setting is moved \textit{verabtim} into the bisupport setting then the resulting partial order is not monotone i.e. composition does not respect the ordering. Call this ordering the support order. Call its dual the cosupport order. These are definable in a bisupport category but in the bisupport category of sets and binary relationships neither of these orderings correspond to relationship inclusion and, worse, composition (of morphisms) is not monotone with respect to these orderings. 

To proceed we need to assume an ordering such that compostion and support and cosupport are monotone. Such a structure we will define as 
\begin{equation}
\label{posetEnrichedBisupportCategory}
\text{poset enriched bisupport category}
\end{equation}
There is a category \cat{pobicat}.

To get to the category of tight zigzags we will need one or more additional conditions.
\workt{I previously thought} that we might need the condition that is currently given as lemma \lref{pointwiseorderingsCompose}:
if 
$
\paralleldiag{a}{b}{f}{f'}
$
\nudgedown{0.6cm}
and
$
\paralleldiag{b}{c}{g}{g'}
$
and $h: a \morph c$ in \catcw 
and if $f' \sleq f$ and $g' \sleq g$ 
and if $h \sleq f \circ g'$ and $h \sleq f' \circ g$ 
then $h \sleq f' \circ g'$. \workt{but this condition does not hold in sets and binary relationships} \workt{The lemma will hold wrt support ordering in a restriction category I think}

\workt{I cannot see this lemma being used anywhere}

\highlight{I need to look at the $h$, $h_L$, $h_R$ lemma and see where this fits in.} 

\workt{Oh no. This doesn't hold in sets and binary relationships. This, I think, is a show stopper.}
\begin{oldtt}
For lack of better terminology when these conditions are met we shall say we have a 
\begin{equation}
\label{wellBehavedPosetEnrichedBisupportCategory}
\text{well-behaved poset enriched bisupport category}
\end{equation}
There is a category \cat{wpobicat}.

Next the category of tight zigzags over a well-behaved poset enriched bisupport 
category has  converses that is compatible with the bisupport structure
\begin{equation}
\label{functionAwareWellbehavedPosetEnrichedBisupportWithConverse}
\text{function aware well-behaved poset enriched bisupport category with compatible converse}
\end{equation}
There is a category \cat{fawpobicatcon}. 

The category of short tight zigzags in addition satisfies condition K.4.
\begin{equation}
\label{functionAwareWellBehavedPosetEnrichedBisupportWithK4Converse}
\text{function-aware well-behaved poset enriched bisupport category with K.4 converse}
\end{equation}
There is a category \cat{fawpobicatk4con}.

The latter is the free 
(\lref{functionAwareWellBehavedPosetEnrichedBisupportWithK4Converse}) 
on a
(\ref{wellBehavedPosetEnrichedBisupportCategory}) 
\end{oldtt}
\end{aside}

\begin{aside}
6th October 2026: We could have co-function aware. K.4 then has a dual. How important is this. Then shortening can be co-shortening of zigzags. What is a 1-1 morphism in a range category? Over a bisupport category does shortening hold for all $f$ that satisfy R.4? Do we therefore absolutely need take account of co-shortening?

Note that I have uncertainties both about the precise positioning of tight zigzags (the place of the $h_L$ and $h_R$ axioma/construction) and about shortening (is cardinality decoration necessary to get a universal construction, what is the remanant of the universal construction when restricted to range categories, is coshortening an absolute requirement(hmmm. Might be a requirment only if we require dual of K.4. Remember K.4 only applies to function-like morphisms.)
\end{aside}
\begin{aside}
I have been working on this since January 2026.
In August 2026 the proofs are fully fleshed out though a few lemmas have proofs trhat are not written up. Through September I am shifting the thinking toward bisupport categories. But then things become less clear.
\begin{enumerate}
  \item lemma \lref{orderingsCompose}  
  The proof that this holds in restriction categories isnt written up. 
  I think I need to do this. 
  Note that $g \leq g'$ implies $f \circ g \leq f \circ g'$ holds only if $f$ is deterministic. A counter example in sets and binary relationships is
  $a = \set{x}, b= \set{y_1,y_2}, c = \set{z_1,z_2}$, 
  $x \mathrel{\overset{f}{\mapsto}} \set{y_1,y_2}$,
  $y_1 \mathrel{\overset{g'}{\mapsto}} \set{z_1}$, 
  $y_2 \mathrel{\overset{g'}{\mapsto}} \set{z_2}$,
  $y_2 \mathrel{\overset{g}{\mapsto}} \set{z_2}$.
  If the ordering is as defined in a restriction category (in which all maps are deterministic as in R.4 holds of them) then whilst $g \leq g'$  it isn't the case that $f \circ g \leq f \circ g'$.
  This needs become an axiom of poset enriched bisupport categories. Proof in
  restriction categories beccomes a proof that restriction categories are poset enriched bisupport catregories.
  \item lemma \lref{pointwiseorderingsCompose}
  I need to prove this for restriction categories. Does it need by an axiom for
  partial order enriched bisupport categroies plus??  
  \item lemma \lref{tighteningCommutesWithComposition} Proof not written up.
  \item lemma \lref{FGsupportLemmas} Proof not written up.
\end{enumerate}
and one definition is missing namely "The Representation of Path Equivalence".

Along the way I had a few expensive false turns including thinking that I could introduce a partial Hom functor. 
This turned out to be a profunctor or something like and so was abandoned.

This paper should probably be split in two with the basic category theory 
of support, cosupport, converse , zigzag construction established in one paper and then used in the other.

Tidying will include
\begin{itemize}
\item Insert missing lemmas.
\begin{itemize}
  \item see lemma tighteningCommutesWithComposition
\end{itemize}
\item Sort out lref and llabel macros. Make them consistent with
indexing of lemmas and uses of lemmas so that I can cross check which lemmas actually used and get rid of unused lemmas.
\begin{itemize}
  \item introduce a implicitref macro
\end{itemize}
\item move zigzag command tests into notation folder.
\item Clean up section regarding category of short tight zigzags as a category with supports and cosupports. Make sure this notion is clear and consistent. Check again with Cockett et al.
\end{itemize}
\end{aside}
\begin{aside}
Much text regarding failed Hom-like functor idea hived off --- see
parallel whenEverythingFellApart.
\end{aside}

\begin{aside}

\begin{enumerate}
\item 
I can add to the definition of bisupport categories a converse operation.
\item the category of zigzags is defined over a bisupport category equipped with a partial order on each hom set.
\item the category of zigzags  is a bisupport category equipped with a converse operation. It might be a free construction of such converse operation.
\item the catgeory of zigzags (as I have so far defined it) relies on the definition of tightness and of tightening. I need tightening because tight zigzags are composed by concatenating them and then by tightening.
\item but why do I actually need tightness in the first place? There is a category of zigzags and a category of tight zigzags. The category of tight zigzags is a quotient of the category of arbitrary zigzags. I know that I am taking this quotient because models in Set will be equivalent but how do I abstract this? Also (later) shortening will have the same 
justification. So is the category of tight zigzags got a position in the hierarchy of real structures or is it just a temporary stopping off point in the development of short tight zigzags. Though shortness is defined wrt range categories it can be introduced for bisupport categories with ordering by only allowing deterministic maps in a shortening. But then jump ahead to short tight zigzags: has this got an abstract reqt. that it is meeting? There is a way of trying to formulate this that has been on the back of my mind for a while.Is it real? 
\item I perhaps need to think of the construction by which I freely generate a converse operation for a category with bisupport.
\end{enumerate}
\end{aside}


\newcommand{\Fquotient}{F/\!\!\sim}  % \! is small negative space

\newcommand{\IfSforCwithRCwords}{If $S$ is a sketch for category \catcw and if $S$ is considered as a data specification with requirement $\reqtc$\ }


\tableofcontents

\begin{aside}
I must user the term zigzag throughout and say that in an RR.5 range category every zigzag has a normal form. These normal forms are the 
short sticklebacks described by Cockett. 

It is potentially confusing though because these normal forms are tuples of morphisms 
of \catcw and these morphisms are generated from the sketch. The morphisms themselves expressed as paths of arrows from the sketch do not in general have normal forms.

\end{aside}

%\include{\srcFolder/MToD.rangeCategories.zigzagcommand.tests.tex}
\include{\srcFolder/MToD.rangeCategories.introduction.tex}
\include{\srcFolder/MToD.rangeCategories.supportCategories.tex}
\include{\srcFolder/MToD.rangeCategories.cosupportCategories.tex}
\include{\srcFolder/MToD.rangeCategories.bisupportCategories.tex}
\include{\srcFolder/MToD.rangeCategories.posetEnrichedBisupportCategories.tex}
\include{\srcFolder/MToD.rangeCategories.zigzags.tex}
\include{\srcFolder/MToD.rangeCategories.bisupportWithConverse.tex}
\include{\srcFolder/MToD.rangeCategories.restrictions_and_support.tex}
\include{\srcFolder/MToD.rangeCategories.ranges_sections_and_weak_range.tex} 
\include{\srcFolder/MToD.rangeCategories.Par_valued_functors.tex}
\include{\srcFolder/MToD.rangeCategories.shortening.tex}

% the file should be renamed the data functor or something more semaantic
\include{\srcFolder/MToD.rangeCategories.carry.tex}
\include{\srcFolder/MToD.rangeCategories.Rel_category.tex}

\section{Goodness of Data Specifications}
\include{\srcFolder/MToD.rangeCategories.goodness_principles.tex}
\include{\srcFolder/MToD.rangeCategories.goodness_criteria.tex}

\section{Goodness Lemmas}
\begin{aside}
Since all these lemmas use the functor $d_a$ then they all have the assumption that
\catcw is an RR.5 range category.
\end{aside}

\include{\srcFolder/MToD.rangeCategories.lemma2A.path_equivalence_representation.tex}

\include{\srcFolder/MToD.rangeCategories.lemma2B.fd_representation.tex}

\include{\srcFolder/MToD.rangeCategories.lemma2C.referential_dependency_representation.tex}

\bibliographystyle{alpha} 
\bibliography{../../../SharedBibliography/temp/bibliography}

\printindex[lemmas]


\end{document}
